\setcounter{secnumdepth}{1} %May be changed to 1 or 2 if section numbers are desired.

\section{Rank-Gap Properties}
\label{app:gap-properties}
Fix a discussion with $N$ candidates and $0<k<N$ selected comments. Assign descending midranks $r_i$ to tied audience scores. The sum of all midranks is $N(N+1)/2$. Midranks can be viewed as expected ranks under uniform permutations within each tied block. For any fixed $k$-set, each such permutation has a rank sum between $k(k+1)/2$ and $k(2N-k+1)/2$; taking expectations preserves these bounds. Subtracting $(k+1)/2$ from the mean selected rank and dividing by $N-k$ therefore gives $0\leq G\leq1$. If the selected set is uniformly random among all $k$-subsets, independently of the fixed ranks, each candidate is included with probability $k/N$. The expected selected mean rank is consequently $(N+1)/2$, giving $\mathbb{E}[G]=1/2$. If all scores tie, every midrank is $(N+1)/2$ and $G=1/2$ for every selection. With no ties, the top and bottom $k$-sets attain the endpoints; ties can prevent this.

\section{Feature Operationalisation and Model Development}
\label{app:feature-operationalisation}
Count and elapsed-time inputs use $\log(1+x)$ transformations. Prior reply composition is the log ratio of earlier replies to earlier roots, with 0.5 added to numerator and denominator. Discussion pace is $\log[(n_{\mathrm{prior}}+0.5)/(h+0.5)]$, where $h$ is hours since publication. Author reception uses the log ratio of retrospectively observed prior upvotes or downvotes to prior author-comment counts, again with 0.5 smoothing. Reply depth is log-transformed and centred among development replies, then multiplied by the reply indicator so that roots have zero centred depth; fitting folds use their own centre.

Lexical diversity and reading difficulty are residualised against log word count using a six-knot cubic spline with ridge regularisation. Missing novelty (i.e. the first comments in time in a discussion) is imputed with the eligible-candidate mean. The twenty AQuA inputs are raw expected ordinal scores on a zero--three scale, not calibrated quality probabilities. Table~\ref{tab:conditional-logit-coefficients-all} lists all 41 scalar predictors and their model-scale units.

\subsection{Development-only sensitivity checks}
\label{app:sensitivity}
All sensitivity checks below use development data only and do not read held-out performance. The Efron--exact comparison uses a reduced model because exact conditional likelihood is computationally expensive: 100 articles sampled from the 207 all-comment articles with $\binom{N_j}{k_j}\leq5000$, with 19 predictors. Efron and exact coefficients had 97.4\% sign agreement, coefficient Spearman correlation 0.995, median absolute coefficient difference 0.010, median query-score Spearman correlation 0.9997, and identical top-scoring comments in 96\% of sampled queries; no sign flip involved a coefficient with absolute value at least 0.10 in either fit.

The collinearity analysis uses all 2,559 development articles and the primary Efron model. For the all-comment scope, the reference-coded full specification showed moderate conditioning: the maximum design condition index was 26.8 (six dimensions above 15 and none above 30). We examined alternative feature and structural specifications excluding author-history or AQuA features, using reply-indicator-only or reply-depth-only terms, and reducing activity controls. These reduced feature-block variants had no fit warnings; removing author-history or reducing activity features caused no coefficient sign changes in the reported selector-specific estimates, while dropping AQuA affected only some small coefficients.

We also evaluated an exactly equivalent midpoint/effect-coded selector parameterisation. Under midpoint coding, the maximum design condition index fell to 11.1, with no dimensions above 15; the robust covariance condition index fell from 14.78 to 9.95. This coding dependence indicates parameterisation geometry rather than a fundamental multicollinearity issue. Selector effect coding leaves the log likelihood unchanged (maximum within-stratum score difference $3.3\times10^{-12}$) and produces identical rankings, so we retain the prespecified full model and report the sensitivity diagnostics rather than treating the reference-coded values as evidence of a failed fit.

\section{Full Conditional-Logit Estimates and Diagnostics}
\label{app:regression}
Coefficients in the following tables describe the development-sample fit; predictive diagnostics refer to the held-out discussions where indicated. Continuous predictors use the shared transformed scale (one development-sample standard deviation); binary indicators compare one with zero. Confidence intervals use article-clustered standard errors. The $p$ columns report two-sided tests and the $q$ columns apply the Benjamini--Hochberg adjustment across all 41 features separately for audience coefficients, editor coefficients, and their differences. The fitted model has finite estimates and standard errors and no recorded fit warnings.

\section{Third-Party Asset Licences}
\label{app:asset-licences}
We record the upstream licence declarations for the frozen language assets used in this study (checked 16 September 2026).
\begin{itemize}
\item \textbf{BGE-M3:} \texttt{BAAI/bge-m3} is released under the MIT licence, as declared in its model card: \url{https://huggingface.co/BAAI/bge-m3}.
\item \textbf{Sentiment:} The CardiffNLP \texttt{twitter-xlm-roberta-base-sentiment} model card links to the XLM-T repository, whose code is licensed under Apache License 2.0: \url{https://github.com/cardiffnlp/xlm-t/blob/main/LICENSE}.
\item \textbf{AQuA:} The upstream repository, which includes the released deliberative-dimension adapters and inference code, carries Apache License 2.0: \url{https://github.com/mabehrendt/AQuA/blob/master/LICENSE}.
\item \textbf{Toxicity:} The TextDetox XLM-R Large Toxicity Classifier v2 model card declares \texttt{openrail++} (OpenRAIL++): \url{https://huggingface.co/textdetox/xlmr-large-toxicity-classifier-v2}.
\end{itemize}

\section{Compute Resources}
\label{app:hardware}
The computational experiments ran on a Linux workstation with an AMD Ryzen 9
9950X processor (16 physical cores, 32 hardware threads) and 64~GB installed
system memory (approximately 60~GiB visible to the operating system). GPU
workloads used one NVIDIA RTX A5000 with 24~GB VRAM.

\onecolumn
\clearpage
\begin{landscape}
\begingroup
\small
% Requires \usepackage{booktabs,longtable,array}

% Standard errors are clustered by article; coefficients are on the

% common conditional-logit log-odds scale.

\setlength{\tabcolsep}{1.5pt}

% ----- all coefficient panel -----

\begin{longtable}{@{}>{\fontsize{6.5}{7.5}\selectfont\raggedright\arraybackslash}p{0.19\linewidth}p{0.04\linewidth}|p{0.15\linewidth}p{0.045\linewidth}p{0.045\linewidth}|p{0.15\linewidth}p{0.045\linewidth}p{0.045\linewidth}|p{0.15\linewidth}p{0.045\linewidth}p{0.045\linewidth}@{}}
\caption{Conditional-logit coefficient estimates} \label{tab:conditional-logit-coefficients-all} \\
\toprule
\multicolumn{2}{c|}{} & \multicolumn{3}{c|}{Journalist} & \multicolumn{3}{c|}{Audience} & \multicolumn{3}{c}{Gap (Journalist-Audience)} \\
Feature & Unit & $\beta_J$ [95\% CI] & p & q & $\beta_A$ [95\% CI] & p & q & $\Delta\beta$ [95\% CI] & p & q \\
\hline
\endfirsthead
\caption[]{Conditional-logit coefficient estimates} \\
\toprule
\multicolumn{2}{c|}{} & \multicolumn{3}{c|}{Journalist} & \multicolumn{3}{c|}{Audience} & \multicolumn{3}{c}{Gap (Journalist-Audience)} \\
Feature & Unit & $\beta_J$ [95\% CI] & p & q & $\beta_A$ [95\% CI] & p & q & $\Delta\beta$ [95\% CI] & p & q \\
\hline
\endhead
\midrule
\multicolumn{11}{r}{Continued on next page} \\
\midrule
\endfoot
\bottomrule
\endlastfoot
\multicolumn{2}{l|}{\textit{Text and language}} & \multicolumn{3}{l|}{} & \multicolumn{3}{l|}{} & \multicolumn{3}{l}{} \\*
Comment length (log words) & 1 SD & 0.590 [0.510, 0.670] & \mbox{$<\!0.001$} & \mbox{$<\!0.001$} & 0.454 [0.375, 0.533] & \mbox{$<\!0.001$} & \mbox{$<\!0.001$} & 0.136 [0.071, 0.201] & \mbox{$<\!0.001$} & \mbox{$<\!0.001$} \\
Positive sentiment & 1 SD & 0.108 [0.076, 0.139] & \mbox{$<\!0.001$} & \mbox{$<\!0.001$} & 0.075 [0.044, 0.107] & \mbox{$<\!0.001$} & \mbox{$<\!0.001$} & 0.032 [0.005, 0.059] & 0.019 & 0.028 \\
Negative sentiment & 1 SD & 0.026 [-0.009, 0.062] & 0.145 & 0.170 & 0.082 [0.047, 0.117] & \mbox{$<\!0.001$} & \mbox{$<\!0.001$} & -0.056 [-0.085, -0.026] & \mbox{$<\!0.001$} & \mbox{$<\!0.001$} \\
Maximum toxicity probability & 1 SD & -0.133 [-0.157, -0.110] & \mbox{$<\!0.001$} & \mbox{$<\!0.001$} & -0.025 [-0.047, -0.003] & 0.025 & 0.035 & -0.108 [-0.127, -0.090] & \mbox{$<\!0.001$} & \mbox{$<\!0.001$} \\
Length-adjusted lexical diversity & 1 SD & 0.006 [-0.010, 0.022] & 0.478 & 0.490 & -0.017 [-0.033, -0.001] & 0.042 & 0.050 & 0.023 [0.010, 0.036] & \mbox{$<\!0.001$} & 0.001 \\
Length-adjusted reading difficulty & 1 SD & 0.042 [0.021, 0.064] & \mbox{$<\!0.001$} & \mbox{$<\!0.001$} & 0.074 [0.052, 0.096] & \mbox{$<\!0.001$} & \mbox{$<\!0.001$} & -0.032 [-0.049, -0.014] & \mbox{$<\!0.001$} & \mbox{$<\!0.001$} \\
URL present & 1 vs 0 & 0.428 [0.302, 0.554] & \mbox{$<\!0.001$} & \mbox{$<\!0.001$} & -0.616 [-0.791, -0.441] & \mbox{$<\!0.001$} & \mbox{$<\!0.001$} & 1.044 [0.892, 1.196] & \mbox{$<\!0.001$} & \mbox{$<\!0.001$} \\
\multicolumn{2}{l|}{\textit{Discussion context and timing}} & \multicolumn{3}{l|}{} & \multicolumn{3}{l|}{} & \multicolumn{3}{l}{} \\*
Article similarity (top-three passages) & 1 SD & 0.121 [0.090, 0.151] & \mbox{$<\!0.001$} & \mbox{$<\!0.001$} & 0.033 [0.002, 0.064] & 0.037 & 0.046 & 0.088 [0.061, 0.114] & \mbox{$<\!0.001$} & \mbox{$<\!0.001$} \\
Novelty from earlier comments & 1 SD & -0.085 [-0.111, -0.059] & \mbox{$<\!0.001$} & \mbox{$<\!0.001$} & -0.124 [-0.149, -0.098] & \mbox{$<\!0.001$} & \mbox{$<\!0.001$} & 0.038 [0.016, 0.060] & \mbox{$<\!0.001$} & 0.001 \\
Hours since article publication & 1 SD & -1.207 [-1.285, -1.130] & \mbox{$<\!0.001$} & \mbox{$<\!0.001$} & -1.541 [-1.629, -1.453] & \mbox{$<\!0.001$} & \mbox{$<\!0.001$} & 0.334 [0.268, 0.399] & \mbox{$<\!0.001$} & \mbox{$<\!0.001$} \\
Earlier comments per article-hour & 1 SD & -0.122 [-0.162, -0.082] & \mbox{$<\!0.001$} & \mbox{$<\!0.001$} & -0.171 [-0.210, -0.133] & \mbox{$<\!0.001$} & \mbox{$<\!0.001$} & 0.050 [0.018, 0.081] & 0.002 & 0.004 \\
Overnight posting period vs workday & 1 vs 0 & 0.170 [0.011, 0.329] & 0.036 & 0.046 & 0.497 [0.348, 0.647] & \mbox{$<\!0.001$} & \mbox{$<\!0.001$} & -0.327 [-0.451, -0.204] & \mbox{$<\!0.001$} & \mbox{$<\!0.001$} \\
Weekday shoulder/evening vs workday & 1 vs 0 & 0.223 [0.130, 0.317] & \mbox{$<\!0.001$} & \mbox{$<\!0.001$} & 0.430 [0.330, 0.529] & \mbox{$<\!0.001$} & \mbox{$<\!0.001$} & -0.207 [-0.284, -0.129] & \mbox{$<\!0.001$} & \mbox{$<\!0.001$} \\
Weekend daytime/evening vs workday & 1 vs 0 & 0.438 [0.087, 0.788] & 0.014 & 0.020 & 0.962 [0.573, 1.351] & \mbox{$<\!0.001$} & \mbox{$<\!0.001$} & -0.524 [-0.809, -0.239] & \mbox{$<\!0.001$} & \mbox{$<\!0.001$} \\
\multicolumn{2}{l|}{\textit{Author history}} & \multicolumn{3}{l|}{} & \multicolumn{3}{l|}{} & \multicolumn{3}{l}{} \\*
Author comments in prior 30 days & 1 SD & -0.074 [-0.095, -0.052] & \mbox{$<\!0.001$} & \mbox{$<\!0.001$} & -0.074 [-0.097, -0.051] & \mbox{$<\!0.001$} & \mbox{$<\!0.001$} & 0.000 [-0.018, 0.019] & 0.975 & 0.975 \\
Upvotes per prior author comment & 1 SD & 0.191 [0.170, 0.211] & \mbox{$<\!0.001$} & \mbox{$<\!0.001$} & 0.259 [0.237, 0.280] & \mbox{$<\!0.001$} & \mbox{$<\!0.001$} & -0.068 [-0.086, -0.050] & \mbox{$<\!0.001$} & \mbox{$<\!0.001$} \\
Downvotes per prior author comment & 1 SD & -0.104 [-0.124, -0.084] & \mbox{$<\!0.001$} & \mbox{$<\!0.001$} & -0.138 [-0.158, -0.118] & \mbox{$<\!0.001$} & \mbox{$<\!0.001$} & 0.034 [0.017, 0.051] & \mbox{$<\!0.001$} & \mbox{$<\!0.001$} \\
Author's earlier comments in article & 1 SD & 0.033 [-0.030, 0.096] & 0.307 & 0.331 & -0.090 [-0.161, -0.019] & 0.013 & 0.020 & 0.123 [0.058, 0.187] & \mbox{$<\!0.001$} & \mbox{$<\!0.001$} \\
\multicolumn{2}{l|}{\textit{Reply structure}} & \multicolumn{3}{l|}{} & \multicolumn{3}{l|}{} & \multicolumn{3}{l}{} \\*
Earlier reply-to-root composition & 1 SD & -0.084 [-0.112, -0.056] & \mbox{$<\!0.001$} & \mbox{$<\!0.001$} & -0.121 [-0.147, -0.094] & \mbox{$<\!0.001$} & \mbox{$<\!0.001$} & 0.037 [0.015, 0.058] & \mbox{$<\!0.001$} & 0.002 \\
Reply rather than root & 1 vs 0 & -4.179 [-4.451, -3.907] & \mbox{$<\!0.001$} & \mbox{$<\!0.001$} & -4.498 [-4.871, -4.125] & \mbox{$<\!0.001$} & \mbox{$<\!0.001$} & 0.319 [-0.063, 0.701] & 0.102 & 0.127 \\
Reply depth centred at the mean & 1 SD & -1.347 [-1.609, -1.084] & \mbox{$<\!0.001$} & \mbox{$<\!0.001$} & -2.248 [-2.608, -1.887] & \mbox{$<\!0.001$} & \mbox{$<\!0.001$} & 0.901 [0.519, 1.282] & \mbox{$<\!0.001$} & \mbox{$<\!0.001$} \\
\multicolumn{2}{l|}{\textit{AQuA dimensions}} & \multicolumn{3}{l|}{} & \multicolumn{3}{l|}{} & \multicolumn{3}{l}{} \\*
Relevance to the discussion topic & 1 SD & 0.040 [-0.010, 0.091] & 0.117 & 0.141 & -0.054 [-0.103, -0.005] & 0.030 & 0.039 & 0.094 [0.051, 0.138] & \mbox{$<\!0.001$} & \mbox{$<\!0.001$} \\
At least one factual claim & 1 SD & -0.027 [-0.105, 0.052] & 0.506 & 0.506 & 0.094 [0.018, 0.170] & 0.016 & 0.022 & -0.121 [-0.189, -0.052] & \mbox{$<\!0.001$} & 0.001 \\
Subjective statement or opinion & 1 SD & 0.159 [0.114, 0.204] & \mbox{$<\!0.001$} & \mbox{$<\!0.001$} & 0.076 [0.033, 0.119] & \mbox{$<\!0.001$} & \mbox{$<\!0.001$} & 0.083 [0.046, 0.120] & \mbox{$<\!0.001$} & \mbox{$<\!0.001$} \\
A statement is justified & 1 SD & 0.163 [0.087, 0.240] & \mbox{$<\!0.001$} & \mbox{$<\!0.001$} & 0.116 [0.045, 0.186] & 0.001 & 0.002 & 0.048 [-0.015, 0.111] & 0.136 & 0.155 \\
A solution is proposed & 1 SD & 0.062 [0.016, 0.109] & 0.009 & 0.013 & -0.036 [-0.080, 0.008] & 0.112 & 0.124 & 0.098 [0.059, 0.137] & \mbox{$<\!0.001$} & \mbox{$<\!0.001$} \\
Additional knowledge is contributed & 1 SD & -0.085 [-0.132, -0.038] & \mbox{$<\!0.001$} & \mbox{$<\!0.001$} & -0.047 [-0.095, 0.001] & 0.054 & 0.062 & -0.038 [-0.078, 0.001] & 0.056 & 0.076 \\
Genuine, non-rhetorical question & 1 SD & -0.066 [-0.095, -0.036] & \mbox{$<\!0.001$} & \mbox{$<\!0.001$} & -0.043 [-0.072, -0.014] & 0.004 & 0.007 & -0.023 [-0.048, 0.002] & 0.075 & 0.097 \\
Reference to another user or community & 1 SD & 0.019 [-0.012, 0.051] & 0.232 & 0.257 & 0.016 [-0.016, 0.047] & 0.332 & 0.359 & 0.004 [-0.024, 0.031] & 0.803 & 0.823 \\
Reference to medium, editor, or moderator & 1 SD & -0.026 [-0.046, -0.006] & 0.010 & 0.014 & -0.003 [-0.023, 0.016] & 0.728 & 0.766 & -0.023 [-0.039, -0.006] & 0.008 & 0.014 \\
Reference to another comment's content & 1 SD & -0.058 [-0.091, -0.025] & \mbox{$<\!0.001$} & \mbox{$<\!0.001$} & -0.079 [-0.112, -0.047] & \mbox{$<\!0.001$} & \mbox{$<\!0.001$} & 0.022 [-0.005, 0.049] & 0.115 & 0.135 \\
Personal reference to another user & 1 SD & -0.018 [-0.046, 0.010] & 0.209 & 0.238 & -0.031 [-0.059, -0.003] & 0.032 & 0.041 & 0.013 [-0.010, 0.036] & 0.270 & 0.291 \\
Reference to tone, spelling, or format & 1 SD & 0.042 [0.004, 0.081] & 0.031 & 0.041 & -0.001 [-0.039, 0.037] & 0.969 & 0.969 & 0.043 [0.010, 0.076] & 0.010 & 0.016 \\
Greeting, farewell, or polite address & 1 SD & 0.045 [0.027, 0.063] & \mbox{$<\!0.001$} & \mbox{$<\!0.001$} & -0.001 [-0.021, 0.018] & 0.902 & 0.924 & 0.046 [0.029, 0.062] & \mbox{$<\!0.001$} & \mbox{$<\!0.001$} \\
Respect, appreciation, or gratitude & 1 SD & 0.046 [0.014, 0.078] & 0.005 & 0.007 & 0.083 [0.051, 0.115] & \mbox{$<\!0.001$} & \mbox{$<\!0.001$} & -0.037 [-0.065, -0.009] & 0.010 & 0.016 \\
Capitalization or punctuation suggesting shouting & 1 SD & -0.046 [-0.066, -0.025] & \mbox{$<\!0.001$} & \mbox{$<\!0.001$} & -0.030 [-0.050, -0.009] & 0.004 & 0.007 & -0.016 [-0.033, 0.001] & 0.069 & 0.092 \\
Language inappropriate for civil discourse & 1 SD & -0.109 [-0.146, -0.072] & \mbox{$<\!0.001$} & \mbox{$<\!0.001$} & -0.085 [-0.120, -0.049] & \mbox{$<\!0.001$} & \mbox{$<\!0.001$} & -0.024 [-0.054, 0.006] & 0.112 & 0.135 \\
Insult or personal attack & 1 SD & 0.021 [-0.024, 0.065] & 0.358 & 0.376 & 0.044 [0.000, 0.088] & 0.048 & 0.056 & -0.023 [-0.060, 0.014] & 0.218 & 0.242 \\
Devaluing, biting mockery & 1 SD & -0.181 [-0.217, -0.145] & \mbox{$<\!0.001$} & \mbox{$<\!0.001$} & -0.144 [-0.179, -0.110] & \mbox{$<\!0.001$} & \mbox{$<\!0.001$} & -0.037 [-0.066, -0.007] & 0.014 & 0.022 \\
Discriminatory language & 1 SD & 0.042 [-0.006, 0.091] & 0.085 & 0.106 & 0.062 [0.015, 0.110] & 0.010 & 0.015 & -0.020 [-0.060, 0.020] & 0.324 & 0.340 \\
Personal story or experience & 1 SD & 0.059 [0.030, 0.089] & \mbox{$<\!0.001$} & \mbox{$<\!0.001$} & 0.088 [0.059, 0.118] & \mbox{$<\!0.001$} & \mbox{$<\!0.001$} & -0.029 [-0.054, -0.004] & 0.021 & 0.030 \\
\end{longtable}

% ----- all fit and diagnostics panel -----

\begin{longtable}{@{}p{0.4\linewidth}p{0.18\linewidth}@{}}
\caption{Conditional-logit diagnostics; robust SEs clustered by article; held-out metrics use test articles} \label{tab:conditional-logit-diagnostics-all} \\
\toprule
Diagnostic & Value \\
\midrule
\endfirsthead
\caption[]{Conditional-logit diagnostics; robust SEs clustered by article; held-out metrics use test articles} \\
\toprule
Diagnostic & Value \\
\midrule
\endhead
\midrule
\multicolumn{2}{r}{Continued on next page} \\
\midrule
\endfoot
\bottomrule
\endlastfoot
Development articles & 2,559 \\
Held-out articles & 2,559 \\
Development candidate comments & 1,189,314 \\
Held-out candidate comments & 1,203,638 \\
Development journalist selections & 10,471 \\
Development audience selections & 10,471 \\
Development selected events & 20,942 \\
Stacked observations & 2,378,628 \\
Stacked selected events & 20,942 \\
Article-selector strata & 5,118 \\
Substantive parameters & 82 \\
Newton-Raphson iterations & 10 \\
Null log likelihood & -123,048.44 \\
Fitted log likelihood & -87,724.12 \\
AIC & 175,612.23 \\
BIC & 176,264.09 \\
Likelihood-ratio chi-square & 70,648.66 (df = 82; $p < 0.001$) \\
Model-based Wald chi-square & 15,849.52 (df = 82; $p < 0.001$) \\
Cluster-robust score chi-square & 2,187.01 (df = 82; $p < 0.001$) \\
Concordance & 0.943 \\
Concordance SE & 0.002 \\
Maximum clustered SE & 0.20 \\
Reference-coded robust correlation-scale condition number & 218.45 \\
Reference-coded maximum robust covariance condition index & 14.78 \\
Midpoint-coded robust correlation-scale condition number & 99.02 \\
Midpoint-coded maximum robust covariance condition index & 9.95 \\
Reference-coded maximum design condition index & 26.77 \\
Midpoint-coded maximum design condition index & 11.09 \\
Reference-coded design dimensions with condition index $\geq 15$ & 6 \\
Midpoint-coded design dimensions with condition index $\geq 15$ & 0 \\
Reference-coded design dimensions with condition index $\geq 30$ & 0 \\
Midpoint-coded design dimensions with condition index $\geq 30$ & 0 \\
Held-out concordance [95\% CI] & 0.942 [0.940, 0.945] \\
Held-out conditional log loss [95\% CI] & 13.46 [13.18, 13.75] \\
\end{longtable}

\endgroup
\end{landscape}

\twocolumn[
\begin{center}
\small
\begin{tabular}{lllll}
\toprule
Model & Audience nDCG@k & Editor nDCG@k & Audience macro-F1 & Editor macro-F1 \\
\midrule
Reg & 0.239 [0.230, 0.249] & 0.233 [0.223, 0.242] & 0.891 [0.888, 0.894] & 0.890 [0.886, 0.892] \\
XGB & 0.358 [0.347, 0.368] & 0.375 [0.365, 0.386] & 0.892 [0.889, 0.895] & 0.897 [0.894, 0.900] \\
XGB-T & \textbf{0.426 [0.415, 0.436]} & \textbf{0.444 [0.433, 0.455]} & 0.895 [0.892, 0.898] & 0.900 [0.897, 0.903] \\
NN & 0.234 [0.225, 0.243] & 0.237 [0.227, 0.247] & 0.898 [0.895, 0.900] & 0.897 [0.894, 0.900] \\
NN-T & 0.263 [0.253, 0.272] & 0.289 [0.279, 0.300] & \textbf{0.912 [0.909, 0.914]} & \textbf{0.913 [0.910, 0.915]} \\
\bottomrule
\end{tabular}

\captionof{table}{Development/training performance. Brackets contain 95\% article-bootstrap intervals. These are fitted-development results, not the cross-validation scores used to choose the ML recipes, and should not be read as estimates of generalisation.}
\label{tab:training-performance}
\end{center}
]

\section{Model Development and Training Performance}
\label{app:development}
The shared article split assigns 2,559 discussions to development and 2,559 to held-out evaluation. Five development folds contain 509, 514, 513, 512, and 511 validation articles. Each fitting fold supplies its own scaling parameters and reply-depth centre. The article split is stratified by month and joint candidate-size tercile and re-randomised until the maximum absolute standardised balance difference is at most 0.05.

For each XGBoost workflow, tuning evaluates 32 reproducible broad random configurations followed by one 27-point local grid centred on the median parameters of the best five broad configurations. The two reported neural winners are selected from a factorial experiment of 64 neural recipes, with reporting restricted to the first audience-label draw. Selection uses development cross-validation, not Table~\ref{tab:training-performance}, which evaluates the final fits on development observations.

The scalar XGBoost winner uses depth 6, 225 trees, learning rate 0.067659, minimum child weight 6.432770, row subsampling 0.712684, column subsampling 0.856588, $L_2$ penalty 13.172306, and zero $L_1$ penalty and split penalty. XGB-T uses depth 5, 291 trees, learning rate 0.061778, minimum child weight 4.742061, row subsampling 0.848434, column subsampling 0.843934, $L_2$ penalty 8.922235, $L_1$ penalty 0.011890, and split penalty 0.456020. Both use histogram trees and the same LambdaRank configuration described in the main text.

Both neural winners use separate selector heads, random sampling of four negatives per positive, and the base network dimensions specified in the main text. Scalar projection is linear--GELU--layer normalisation--dropout; selector heads are linear--GELU--dropout--linear. The plateau scheduler halves the learning rate after two validation epochs without improvement, down to $10^{-6}$. Development folds allow at most 30 epochs with early-stopping patience eight; the final full-development fit uses the median selected epoch across folds (NN: nine; NN-T: twelve).

The development-CV mean nDCG, averaging selectors, is 0.24495 for XGB, 0.25167 for XGB-T, 0.23400 for NN, and 0.23935 for NN-T. These validation scores are distinct from the fitted-development scores. The latter are particularly higher for the boosted trees and should not be used to infer predictive performance on new discussions.

The fitted-development table (Table~\ref{tab:training-performance}) makes the generalisation gap visible. XGB-T has fitted-development nDCG@$k$ of 0.444 for editors and 0.426 for audiences, compared with held-out values of 0.257 and 0.240. The regression's corresponding development and held-out values are 0.233/0.239 for editors/audiences and 0.232/0.232, respectively. Model flexibility therefore improves fitted development performance much more than recovery in held-out discussions; the two tables use the same complete-discussion nDCG diagnostic, while the macro-F1 columns use the separate balanced diagnostic.
